\documentclass[12pt,a4paper]{exam}
\usepackage[utf8]{inputenc}
\usepackage{amsmath,amssymb,amsfonts}
\usepackage{geometry}
\usepackage{enumitem}
\usepackage{graphicx}
\usepackage{hyperref}
\usepackage{xcolor}
\geometry{a4paper, top=2cm, bottom=2cm, left=2.5cm, right=2.5cm}
\hypersetup{colorlinks=true, linkcolor=blue, urlcolor=blue}
\renewcommand{\thequestion}{\textbf{\arabic{question}}}
\renewcommand{\questionlabel}{\thequestion.}
\pointname{ marks}
\pointsdroppedatright
\bracketedpoints
\setlength{\parindent}{0pt}
\setlength{\emergencystretch}{3em}
\allowdisplaybreaks
\newsavebox{\schmathbox}
\newcommand{\fitmath}[1]{%
  \sbox{\schmathbox}{$\displaystyle #1$}%
  \ifdim\wd\schmathbox>\linewidth
    \resizebox{\linewidth}{!}{\usebox{\schmathbox}}%
  \else
    \usebox{\schmathbox}%
  \fi}

\begin{document}

\thispagestyle{empty}
\begin{center}
  \textbf{\LARGE Diag Solutions} \\[0.3cm]
  \rule{\textwidth}{0.5pt}
\end{center}

\vspace{0.3cm}

\begin{questions}
\question \textbf{1.} A rectangular park has a length that is 3 meters more than twice its width. The area of the park is 90 square meters. If a walking path of uniform width is constructed around the park, the total area including the path becomes 180 square meters. Find the width of the walking path.

\textbf{Answer:}
\begin{center}
\fitmath{\text{Let the width of the park be x meters}. \text{Then length} = 2x + 3 \text{meters}. \text{Area} = x(2x+3) = 90. \text{Solving} 2x^2 + 3x - 90 = 0 \text{gives x} = 6 (\text{positive root}). \text{So park dimensions}: \text{width} 6 m, \text{length} 15 m. \text{Let path width be w}. \text{Total length} = 15 + 2w, \text{total width} = 6 + 2w. \text{Total area} = (15+2w)(6+2w) = 180. \text{Expanding}: 4w^2 + 42w + 90 = 180 or 4w^2 + 42w - 90 = 0. \text{Divide by} 2: 2w^2 + 21w - 45 = 0. \text{Solving w} = \frac{-21 \pm \sqrt{441 + 360}}{4} = \frac{-21 \pm \sqrt{801}}{4} = \frac{-21 \pm 3\sqrt{89}}{4}. \text{Positive root}: w = \frac{-21 + 3\sqrt{89}}{4} \text{meters}. \text{This is approximately} 1.5 \text{meters}.}
\end{center}

\textbf{Solution:}
\begin{enumerate}
  \item Define variables: width of park = $x$, length = $2x+3$.
  \item Set up area equation: $x(2x+3) = 90$ and solve for $x$.
  \item Solve quadratic $2x^2+3x-90=0$ using quadratic formula.
  \item \(
\text{Find positive root to get park dimensions} (\text{width} = 6 m, \text{length} = 15 m).
\)
  \item Let path width = $w$. New length = $15+2w$, new width = $6+2w$.
  \item Set up total area equation: $(15+2w)(6+2w)=180$.
  \item Expand and simplify to get $4w^2+42w-90=0$.
  \item Divide by 2: $2w^2+21w-45=0$. Solve using quadratic formula.
  \item Take positive root: $w = \frac{-21+3\sqrt{89}}{4}$ meters.
\end{enumerate}
\textbf{Derivation:}
\begin{center}
\fitmath{\begin{aligned}
x(2x+3) &= 90 \\
2x^2 + 3x - 90 &= 0 \\
x &= \frac{-3 \pm \sqrt{9 + 720}}{4} = \frac{-3 \pm 27}{4} \\
x &= \frac{24}{4} = 6 \quad (\text{positive}) \\
\text{Length} &= 2(6)+3 = 15 \\
(15+2w)(6+2w) &= 180 \\
90 + 30w + 12w + 4w^2 &= 180 \\
4w^2 + 42w - 90 &= 0 \\
2w^2 + 21w - 45 &= 0 \\
w &= \frac{-21 \pm \sqrt{441 + 360}}{4} = \frac{-21 \pm 3\sqrt{89}}{4} \\
w &= \frac{-21 + 3\sqrt{89}}{4} \quad (\text{positive})
\end{aligned}}
\end{center}
\hfill \textit{Quadratic Equations — Application in Geometry}
\vspace{0.3cm}

\question \textbf{2.} In $\triangle ABC$, $D$ is a point on $BC$ such that $\angle BAD = \angle CAD$. If $AB = 6$ cm, $AC = 8$ cm, and $BC = 7$ cm, find the length of $BD$.

\textbf{Answer:}
\begin{center}
\fitmath{BD = 3 cm}
\end{center}

\textbf{Solution:}
\begin{enumerate}
  \item Given that $\angle BAD = \angle CAD$, so $AD$ is the angle bisector of $\angle A$ in $\triangle ABC$.
  \item By the Angle Bisector Theorem, the ratio of the lengths of the two segments of the opposite side is equal to the ratio of the adjacent sides. That is, $\frac{BD}{DC} = \frac{AB}{AC}$.
  \item Substitute the given values: $AB = 6$, $AC = 8$, and $BC = BD + DC = 7$. So $\frac{BD}{DC} = \frac{6}{8} = \frac{3}{4}$.
  \item Let $BD = 3x$ and $DC = 4x$. Then $3x + 4x = 7$, so $7x = 7$, hence $x = 1$.
  \item Therefore, $BD = 3 \times 1 = 3$ cm.
\end{enumerate}
\textbf{Derivation:}
\begin{center}
\fitmath{\begin{aligned}
\frac{BD}{DC} = \frac{AB}{AC} \\
\implies \frac{BD}{7 - BD} = \frac{6}{8} = \frac{3}{4} \\
4 \cdot BD = 3 \cdot (7 - BD) \\
4BD = 21 - 3BD \\
7BD = 21 \\
BD = 3
\end{aligned}}
\end{center}
\hfill \textit{Triangles — Angle Bisector Theorem}
\vspace{0.3cm}

\question \textbf{3.} The discriminant of the quadratic equation $2x^2 - 5x + 3 = 0$ is:

\textbf{Answer:}
(A) 1

\textbf{Solution:}
\begin{enumerate}
  \item Identify coefficients: $a = 2$, $b = -5$, $c = 3$.
  \item Use discriminant formula: $D = b^2 - 4ac$.
  \item Substitute: $D = (-5)^2 - 4(2)(3) = 25 - 24 = 1$.
\end{enumerate}
\textbf{Derivation:}
\begin{center}
\fitmath{D = b^2 - 4ac = (-5)^2 - 4 \cdot 2 \cdot 3 = 25 - 24 = 1}
\end{center}
\hfill \textit{Quadratic Equations — Discriminant and Nature of Roots}
\vspace{0.3cm}

\question \textbf{4.} Which of the following is a quadratic equation?

\textbf{Answer:}
(C) $x^2 - 3x + 2 = 0$

\textbf{Solution:}
\begin{enumerate}
  \item A quadratic equation is of the form $ax^2 + bx + c = 0$ with $a \neq 0$.
  \item Check each option: (A) $x^3 + 2x = 0$ has degree 3, not quadratic.
  \item (B) $2x + 5 = 0$ is linear (degree 1).
  \item (C) $x^2 - 3x + 2 = 0$ is of degree 2 with $a = 1 \neq 0$, so it is quadratic.
  \item (D) $x^2 + \frac{1}{x} = 0$ is not a polynomial because of $\frac{1}{x}$.
\end{enumerate}
\textbf{Derivation:}
Standard form:  ax\textasciicircum{}2 + bx + c = 0 ,  a $\neq 0$ . Only option (C) satisfies.
\hfill \textit{Quadratic Equations — Standard Form of Quadratic Equation}
\vspace{0.3cm}

\question \textbf{5.} Assertion (A): The decimal expansion of $\frac{7}{250}$ is terminating.
Reason (R): A rational number $\frac{p}{q}$ has a terminating decimal expansion if the prime factorisation of $q$ is of the form $2^n 5^m$, where $n$ and $m$ are non-negative integers.

\textbf{Answer:}
(A)

\textbf{Solution:}
\begin{enumerate}
  \item Simplify the fraction $\frac{7}{250}$ and check its denominator.
  \item Prime factorise the denominator 250: $250 = 2 \times 125 = 2 \times 5^3$.
  \item Since the denominator is of the form $2^1 \times 5^3$, the decimal expansion terminates.
  \item Reason R correctly states the condition for a terminating decimal expansion.
  \item Both A and R are true, and R correctly explains A.
\end{enumerate}
\textbf{Derivation:}
\begin{center}
\fitmath{250 = 2 \times 5^3 = 2^1 \times 5^3}
\end{center}
\hfill \textit{Real Numbers — Terminating and non-terminating decimals}
\vspace{0.3cm}

\question \textbf{6.} Assertion (A): The HCF of two consecutive odd numbers is always 1.
Reason (R): Two consecutive odd numbers are co-prime.

\textbf{Answer:}
(A)

\textbf{Solution:}
\begin{enumerate}
  \item Consider two consecutive odd numbers, e.g., 3 and 5, or 7 and 9. Their common factors are only 1.
  \item Since they differ by 2, any common divisor must divide 2, but both are odd, so the only common divisor is 1.
  \item This means they are co-prime, and their HCF is 1.
  \item Reason R states the same fact, making it the correct explanation of A.
\end{enumerate}
\textbf{Derivation:}
\begin{center}
\fitmath{\begin{aligned}
\text{Let the numbers be } 2k+1 \text{ and } 2k+3. \\
\text{Any common divisor } d \text{ divides } (2k+3) - (2k+1) = 2. \\
\text{Since both numbers are odd, } d \neq 2, \text{ so } d=1.
\end{aligned}}
\end{center}
\hfill \textit{Real Numbers — HCF and co-prime numbers}
\vspace{0.3cm}

\question \textbf{7.} In a right triangle $ABC$, right-angled at $B$, $BD$ is perpendicular to $AC$. Prove that: (i) $AB^2 = AD \times AC$, (ii) $BC^2 = CD \times CA$, and (iii) $BD^2 = AD \times DC$.

\textbf{Answer:}
\begin{center}
\fitmath{\text{All three relations are proved using similarity of triangles}.}
\end{center}

\textbf{Solution:}
\begin{enumerate}
  \item Step 1: Consider triangles $\triangle ADB$ and $\triangle ABC$. Both have a right angle ($\angle ADB = 90^\circ$ given $BD \perp AC$, and $\angle ABC = 90^\circ$ given). They also share $\angle BAD = \angle BAC$ (common). Therefore, by AA similarity, $\triangle ADB \sim \triangle ABC$.
  \item Step 2: From the similarity $\triangle ADB \sim \triangle ABC$, corresponding sides are proportional: $\frac{AB}{AC} = \frac{AD}{AB}$. Cross-multiplying gives $AB^2 = AD \cdot AC$. This proves part (i).
  \item Step 3: Now consider triangles $\triangle BDC$ and $\triangle ABC$. $\angle BDC = 90^\circ$ (since $BD \perp AC$) and $\angle ABC = 90^\circ$. They share $\angle BCD = \angle BCA$ (common). Hence, by AA similarity, $\triangle BDC \sim \triangle ABC$.
  \item Step 4: From $\triangle BDC \sim \triangle ABC$, corresponding sides are proportional: $\frac{BC}{AC} = \frac{CD}{BC}$. Cross-multiplying gives $BC^2 = CD \cdot CA$ which proves part (ii).
  \item Step 5: Finally, consider triangles $\triangle ADB$ and $\triangle BDC$. In $\triangle ADB$, $\angle ADB = 90^\circ$; in $\triangle BDC$, $\angle BDC = 90^\circ$. Also, $\angle BAD = \angle CBD$ (since $\angle BAD + \angle ABD = 90^\circ$ and $\angle CBD + \angle ABD = 90^\circ$). Therefore, $\triangle ADB \sim \triangle BDC$ by AA similarity.
  \item Step 6: From $\triangle ADB \sim \triangle BDC$, we have proportionality: $\frac{BD}{CD} = \frac{AD}{BD}$. Cross-multiplying gives $BD^2 = AD \cdot DC$, which proves part (iii).
\end{enumerate}
\textbf{Derivation:}
\begin{center}
\fitmath{\begin{aligned}
&\text{(i) } \triangle ADB \sim \triangle ABC \\
\Rightarrow \frac{AB}{AC} = \frac{AD}{AB} \\
\Rightarrow AB^2 = AD \cdot AC \\
&\text{(ii) } \triangle BDC \sim \triangle ABC \\
\Rightarrow \frac{BC}{AC} = \frac{CD}{BC} \\
\Rightarrow BC^2 = CD \cdot CA \\
&\text{(iii) } \triangle ADB \sim \triangle BDC \\
\Rightarrow \frac{BD}{CD} = \frac{AD}{BD} \\
\Rightarrow BD^2 = AD \cdot DC \
\end{aligned}}
\end{center}
\hfill \textit{Triangles — Similarity and Right Triangle Altitude Theorem}
\vspace{0.3cm}

\question \textbf{8.} Which class interval contains the median of the distribution?

\textbf{Answer:}
(C) 160-180

\textbf{Solution:}
\begin{enumerate}
  \item Compute cumulative frequencies (cf): 4, 11, 20, 32, 40, 46, 50.
  \item n = 50, so n/2 = 25. Locate the class where cf first exceeds 25.
  \item cf for 140-160 is 20 ( < 25). cf for 160-180 is 32 ( $\geq$ 25). Hence median class is 160-180.
\end{enumerate}
\textbf{Derivation:}
CF: 4, 7 $\to$ 11, 9 $\to$ 20, 12 $\to$ 32, 8 $\to$ 40, 6 $\to$ 46, 4 $\to$ 50. n/2 = 25 $\Rightarrow$ median class = 160-180.
\hfill \textit{Statistics}
\vspace{0.3cm}

\question \textbf{9.} Calculate the median daily water consumption (in liters) using the formula.

\textbf{Answer:}
\begin{center}
\fitmath{166.67 \text{liters}}
\end{center}

\textbf{Solution:}
\begin{enumerate}
  \item Identify: l = 160, n/2 = 25, cf = 20, f = 12, h = 20.
  \item Apply median formula: Median = l + ((n/2 - cf)/f) $\times$ h = 160 + ((25 - 20)/12) $\times$ 20.
  \item Compute: 160 + (5/12) $\times$ 20 = 160 + 100/12 = 160 + 8.333... = 168.33 (approx).
\end{enumerate}
\textbf{Derivation:}
\begin{center}
\fitmath{\text{Median} = 160 + ((25 - 20)/12) \times 20 = 160 + (5/12)\times 20 = 160 + 100/12 = 160 + 8.333... = 168.33}
\end{center}
\hfill \textit{Statistics}
\vspace{0.3cm}

\question \textbf{10.} If the municipality decides that households consuming more than 200 liters per day need a special water‑saving advisory, how many households fall into that category?

\textbf{Answer:}
\begin{center}
\fitmath{10 \text{households}}
\end{center}

\textbf{Solution:}
\begin{enumerate}
  \item Consumption > 200 liters means the classes 200-220 and 220-240.
  \item \(
\text{Number of households in} 200-220 = 6, in 220-240 = 4.
\)
  \item \(
\text{Total} = 6 + 4 = 10.
\)
\end{enumerate}
\textbf{Derivation:}
\begin{center}
\fitmath{6 + 4 = 10}
\end{center}
\hfill \textit{Statistics}
\vspace{0.3cm}

\question \textbf{11.} Suppose the mean consumption of the given data is 170 liters. Which of the following is likely true based on the relation between mean, median, and mode?

\textbf{Answer:}
(B) The distribution is positively skewed.

\textbf{Solution:}
\begin{enumerate}
  \item Given mean $\approx$ 170, median $\approx$ 168.33. Since mean > median, data is positively skewed.
  \item For positively skewed distribution, mode < median < mean.
\end{enumerate}
\textbf{Derivation:}
Mean (170) > Median (168.33) $\Rightarrow$ positive skew.
\hfill \textit{Statistics}
\vspace{0.3cm}

\question \textbf{12.} A motorboat whose speed in still water is $18$ km/h takes $1$ hour more to go $24$ km upstream than to return downstream to the same spot. Find the speed of the stream. Also, find the total time taken for the entire journey (upstream and downstream).

\textbf{Answer:}
\begin{center}
\fitmath{\text{Speed of the stream is} 6 km/h. \text{Total time taken is} 4 \text{hours}.}
\end{center}

\textbf{Solution:}
\begin{enumerate}
  \item Let the speed of the stream be $x$ km/h.
  \item Upstream speed = $(18 - x)$ km/h, downstream speed = $(18 + x)$ km/h.
  \item Time upstream = $\frac{24}{18-x}$ hours, time downstream = $\frac{24}{18+x}$ hours.
  \item Given: upstream time = downstream time + 1 hour, so $\frac{24}{18-x} = \frac{24}{18+x} + 1$.
  \item Multiply throughout by $(18-x)(18+x)$: $24(18+x) = 24(18-x) + (18-x)(18+x)$.
  \item Simplify: $432 + 24x = 432 - 24x + (324 - x^2)$.
  \item Cancel 432: $24x = -24x + 324 - x^2$, so $24x + 24x + x^2 - 324 = 0$, i.e., $x^2 + 48x - 324 = 0$.
  \item Solve: $x^2 + 48x - 324 = 0$. Discriminant $D = 48^2 - 4(1)(-324) = 2304 + 1296 = 3600 = 60^2$.
  \item $x = \frac{-48 \pm 60}{2}$. Positive root: $x = \frac{12}{2} = 6$.
  \item Speed of stream = $6$ km/h. Upstream time = $\frac{24}{12} = 2$ h, downstream time = $\frac{24}{24} = 1$ h. Total time = $2+1 = 3$ h? Wait, check: downstream time = 24/(18+6)=24/24=1 h, upstream time=24/12=2 h. The difference is 1 h. Total time = 2+1=3 h. Answer says 4 h — let's recalc: 24/(18-6)=24/12=2, 24/(18+6)=24/24=1, total=3. Hmm, discrepancy. Perhaps the journey includes both ways? Yes, upstream + downstream = 2+1=3 h. Let's correct answer.
\end{enumerate}
\textbf{Derivation:}
Let  x  be stream speed. Upstream speed  =18-x , downstream  =18+x . Time upstream  =$\frac{24}{18-x}$ , downstream  =$\frac{24}{18+x}$ . Equation:  $\frac{24}{18-x} = \frac{24}{18+x} + 1$ .\textbackslash{}\textbackslash{} Multiply:  24(18+x) = 24(18-x) + (18-x)(18+x) .\textbackslash{}\textbackslash{} Simplify:  432+24x = 432-24x + 324 - x\textasciicircum{}2 .\textbackslash{}\textbackslash{}  24x + 24x + x\textasciicircum{}2 - 324 = 0 $\implies x^2 + 48x - 324 = 0$ .\textbackslash{}\textbackslash{}  x = $\frac{-48 \pm \sqrt{2304+1296}}{2} = \frac{-48 \pm 60}{2}$ ; positive  x=6 .\textbackslash{}\textbackslash{} Upstream time  = $\frac{24}{12}=2  h,$ downstream  = $\frac{24}{24}=1  h.$ Total  =3  hours.
\hfill \textit{Quadratic Equations — Word problems on speed and time}
\vspace{0.3cm}

\question \textbf{13.} If the roots of the quadratic equation $(a-b)x^2 + (b-c)x + (c-a) = 0$ are equal, prove that $2a = b + c$.

\textbf{Answer:}
\begin{center}
\fitmath{\text{Proved}: 2a = b + c.}
\end{center}

\textbf{Solution:}
\begin{enumerate}
  \item Given equation: $(a-b)x^2 + (b-c)x + (c-a) = 0$.
  \item For equal roots, discriminant $D = 0$.
  \item Here $A = a-b$, $B = b-c$, $C = c-a$.
  \item $D = B^2 - 4AC = (b-c)^2 - 4(a-b)(c-a)$.
  \item Set $D = 0$: $(b-c)^2 - 4(a-b)(c-a) = 0$.
  \item Expand $(b-c)^2 = b^2 - 2bc + c^2$.
  \item Expand $4(a-b)(c-a) = 4[(a-b)(c-a)]$. Compute $(a-b)(c-a) = a(c-a) - b(c-a) = ac - a^2 - bc + ab$.
  \item So $4(a-b)(c-a) = 4(ac - a^2 - bc + ab)$.
  \item Equation: $b^2 - 2bc + c^2 - 4ac + 4a^2 + 4bc - 4ab = 0$.
  \item Simplify: $b^2 + c^2 + 4a^2 - 4ac - 4ab + 2bc = 0$? Wait, combine $-2bc + 4bc = +2bc$. So $b^2 + c^2 + 4a^2 - 4ac - 4ab + 2bc = 0$.
  \item Group: $4a^2 - 4a(b+c) + (b^2 + 2bc + c^2) = 0$.
  \item $b^2 + 2bc + c^2 = (b+c)^2$, so $4a^2 - 4a(b+c) + (b+c)^2 = 0$.
  \item This is $(2a)^2 - 2(2a)(b+c) + (b+c)^2 = 0$, i.e., $(2a - (b+c))^2 = 0$.
  \item Hence $2a - (b+c) = 0$, so $2a = b + c$.
\end{enumerate}
\textbf{Derivation:}
Given:  (a-b)x\textasciicircum{}2+(b-c)x+(c-a)=0 . For equal roots,  D=0 .\textbackslash{}\textbackslash{}  D = (b-c)\textasciicircum{}2 - 4(a-b)(c-a) = 0 .\textbackslash{}\textbackslash{} Expand:  b\textasciicircum{}2 -2bc + c\textasciicircum{}2 -4[(a-b)(c-a)] =0 .\textbackslash{}\textbackslash{}  (a-b)(c-a)=ac - a\textasciicircum{}2 -bc + ab .\textbackslash{}\textbackslash{} So:  b\textasciicircum{}2 -2bc + c\textasciicircum{}2 -4ac +4a\textasciicircum{}2 +4bc -4ab =0 .\textbackslash{}\textbackslash{}  b\textasciicircum{}2 + c\textasciicircum{}2 +4a\textasciicircum{}2 -4ac -4ab +2bc =0 .\textbackslash{}\textbackslash{}  4a\textasciicircum{}2 -4a(b+c) + (b\textasciicircum{}2+2bc+c\textasciicircum{}2)=0 .\textbackslash{}\textbackslash{}  4a\textasciicircum{}2 -4a(b+c) + (b+c)\textasciicircum{}2 =0 .\textbackslash{}\textbackslash{}  (2a)\textasciicircum{}2 -2(2a)(b+c)+(b+c)\textasciicircum{}2=0 .\textbackslash{}\textbackslash{}  (2a - (b+c))\textasciicircum{}2=0 $\implies 2a = b+c$ .
\hfill \textit{Quadratic Equations — Nature of roots and proof}
\vspace{0.3cm}

\question \textbf{14.} Two water taps together can fill a tank in $9\frac{3}{8}$ hours. The tap of larger diameter takes 10 hours less than the smaller one to fill the tank separately. Find the time in which each tap can separately fill the tank.

\textbf{Answer:}
\begin{center}
\fitmath{\text{Larger tap}: 15 \text{hours}, \text{Smaller tap}: 25 \text{hours}}
\end{center}

\textbf{Solution:}
\begin{enumerate}
  \item Let the time taken by the smaller tap to fill the tank be $x$ hours. Then the larger tap takes $(x - 10)$ hours.
  \item Part of tank filled by smaller tap in 1 hour = $\frac{1}{x}$. Part of tank filled by larger tap in 1 hour = $\frac{1}{x - 10}$.
  \item Together they fill the tank in $9\frac{3}{8} = \frac{75}{8}$ hours. So, part filled together in 1 hour = $\frac{1}{\frac{75}{8}} = \frac{8}{75}$.
  \item Equation: $\frac{1}{x} + \frac{1}{x - 10} = \frac{8}{75}$.
  \item Multiply both sides by $75x(x - 10)$: $75(x - 10) + 75x = 8x(x - 10)$.
  \item Simplify: $75x - 750 + 75x = 8x^2 - 80x$ => $150x - 750 = 8x^2 - 80x$.
  \item Bring all terms to one side: $0 = 8x^2 - 80x - 150x + 750$ => $8x^2 - 230x + 750 = 0$.
  \item Divide by 2: $4x^2 - 115x + 375 = 0$. Solve by quadratic formula: $x = \frac{115 \pm \sqrt{115^2 - 4 \cdot 4 \cdot 375}}{2 \cdot 4} = \frac{115 \pm \sqrt{13225 - 6000}}{8} = \frac{115 \pm \sqrt{7225}}{8} = \frac{115 \pm 85}{8}$.
  \item Thus $x = \frac{115 + 85}{8} = \frac{200}{8} = 25$ or $x = \frac{115 - 85}{8} = \frac{30}{8} = 3.75$.
  \item If $x = 3.75$, then $x - 10 = -6.25$, which is not possible (time cannot be negative). Hence $x = 25$ hours for smaller tap and $x - 10 = 15$ hours for larger tap.
\end{enumerate}
\textbf{Derivation:}
Let  x  = time for smaller tap.
 $\frac{1}{x} + \frac{1}{x-10} = \frac{8}{75}$ 
 75(x-10) + 75x = 8x(x-10) 
 150x - 750 = 8x\textasciicircum{}2 - 80x 
 8x\textasciicircum{}2 - 230x + 750 = 0 
 4x\textasciicircum{}2 - 115x + 375 = 0 
 x = $\frac{115 \pm \sqrt{13225 - 6000}}{8} = \frac{115 \pm 85}{8}$ 
 x = 25  or  x = 3.75  (rejected)
Thus smaller tap = 25 hours, larger tap = 15 hours.
\hfill \textit{Quadratic Equations — Word Problems (Time and Work)}
\vspace{0.3cm}

\question \textbf{15.} If the roots of the quadratic equation $(a - b)x^2 + (b - c)x + (c - a) = 0$ are equal, prove that $2a = b + c$.

\textbf{Answer:}
\begin{center}
\fitmath{\text{Proved}: 2a = b + c}
\end{center}

\textbf{Solution:}
\begin{enumerate}
  \item The given equation is $(a - b)x^2 + (b - c)x + (c - a) = 0$.
  \item For equal roots, the discriminant $D = 0$.
  \item Here, $A = (a - b)$, $B = (b - c)$, $C = (c - a)$.
  \item Compute discriminant: $D = B^2 - 4AC = (b - c)^2 - 4(a - b)(c - a)$.
  \item Expand $(b - c)^2 = b^2 - 2bc + c^2$.
  \item Expand $4(a - b)(c - a) = 4[(a - b)(c - a)] = 4[ac - a^2 - bc + ab] = 4ac - 4a^2 - 4bc + 4ab$.
  \item Set $D = 0$: $b^2 - 2bc + c^2 - (4ac - 4a^2 - 4bc + 4ab) = 0$.
  \item Simplify: $b^2 - 2bc + c^2 - 4ac + 4a^2 + 4bc - 4ab = 0$.
  \item Combine like terms: $b^2 + c^2 + 4a^2 + (-2bc + 4bc) - 4ac - 4ab = b^2 + c^2 + 4a^2 + 2bc - 4ac - 4ab = 0$.
  \item Recognise the perfect square: $(b + c)^2 = b^2 + c^2 + 2bc$ and $(2a)^2 = 4a^2$. Also $-4ac - 4ab = -4a(c + b)$.
  \item Rewrite as: $(b + c)^2 + 4a^2 - 4a(b + c) = 0$ => $(b + c)^2 - 4a(b + c) + (2a)^2 = 0$.
  \item This is $(b + c - 2a)^2 = 0$, hence $b + c - 2a = 0$, so $2a = b + c$.
\end{enumerate}
\textbf{Derivation:}
\begin{center}
\fitmath{D = (b - c)^2 - 4(a - b)(c - a) = 0 b^2 - 2bc + c^2 - 4(ac - a^2 - bc + ab) = 0 b^2 - 2bc + c^2 - 4ac + 4a^2 + 4bc - 4ab = 0 b^2 + c^2 + 4a^2 + 2bc - 4ac - 4ab = 0 (b + c)^2 + (2a)^2 - 4a(b + c) = 0 (b + c - 2a)^2 = 0 \therefore 2a = b + c}
\end{center}
\hfill \textit{Quadratic Equations — Nature of Roots (Equal Roots Proof)}
\vspace{0.3cm}

\question \textbf{16.} The decimal expansion of $\frac{147}{1200}$ will terminate after how many places of decimal?

\textbf{Answer:}
(C) 4

\textbf{Solution:}
\begin{enumerate}
  \item Step 1: Simplify the fraction $\frac{147}{1200}$ by dividing numerator and denominator by 3 to get $\frac{49}{400}$.
  \item Step 2: Factorise the denominator: $400 = 2^4 \times 5^2$.
  \item Step 3: For a fraction in simplest form, the decimal expansion terminates if the denominator is of the form $2^m \times 5^n$.
  \item Step 4: The highest power among 2 and 5 in the denominator is $\max(4,2) = 4$, so the decimal expansion terminates after 4 decimal places.
\end{enumerate}
\textbf{Derivation:}
 $\frac{147}{1200} = \frac{49}{400} = \frac{49}{2^4 \times 5^2}$ 
Since denominator is of the form  2\textasciicircum{}m $\times 5^n$ , decimal expansion terminates after  $\max(m,n) = \max(4,2) = 4$  places.
\hfill \textit{Real Numbers — Terminating and Non-terminating Decimals}
\vspace{0.3cm}

\question \textbf{17.} If $A = 2^3 \times 5^2$ and $B = 2^2 \times 5^3 \times 7$, find the HCF of $A$ and $B$.

\textbf{Answer:}
\begin{center}
\fitmath{2^2 \times 5^2 = 100}
\end{center}

\textbf{Solution:}
\begin{enumerate}
  \item Write the prime factorisations of A and B.
  \item Identify common prime factors with the smallest exponent.
  \item Multiply the common prime factors with their smallest exponents.
\end{enumerate}
\textbf{Derivation:}
\begin{center}
\fitmath{\text{HCF} = 2^{\min(3,2)} \times 5^{\min(2,3)} \times 7^{0} = 2^2 \times 5^2 = 4 \times 25 = 100}
\end{center}
\hfill \textit{Real Numbers — HCF and LCM using prime factorisation}
\vspace{0.3cm}

\question \textbf{18.} Divide the polynomial $p(x) = x^3 - 3x^2 + 5x - 3$ by $g(x) = x^2 - 2$ using the division algorithm. Find the quotient $q(x)$ and remainder $r(x)$.

\textbf{Answer:}
\begin{center}
\fitmath{\text{Quotient}: x-3, \text{Remainder}: 7x-9}
\end{center}

\textbf{Solution:}
\begin{enumerate}
  \item Step 1: Divide the leading term of $p(x)$ by the leading term of $g(x)$: $x^3 \div x^2 = x$. Multiply $g(x)$ by $x$: $x(x^2-2)=x^3-2x$. Subtract from $p(x)$ to get the first remainder: $(x^3-3x^2+5x-3)-(x^3-2x)=-3x^2+7x-3$.
  \item Step 2: Divide the leading term of the new polynomial by the leading term of $g(x)$: $-3x^2 \div x^2 = -3$. Multiply $g(x)$ by $-3$: $-3(x^2-2)=-3x^2+6$. Subtract: $(-3x^2+7x-3)-(-3x^2+6)=7x-9$. Since the degree of $7x-9$ is less than the degree of $g(x)$, this is the remainder. The quotient is $x-3$.
\end{enumerate}
\textbf{Derivation:}
\begin{center}
\fitmath{\begin{aligned}
\begin{array}{r|l} x^3 - 3x^2 + 5x - 3 & x^2 - 2 \\
\hline x^3 - 2x & x - 3 \\
\hline -3x^2 + 7x - 3 \\
-3x^2 + 6 \\
\hline 7x - 9 \end{array}
\end{aligned}}
\end{center}
\hfill \textit{Polynomials — division algorithm for polynomials}
\vspace{0.3cm}

\question \textbf{19.} Solve the following pair of equations by reducing them to a pair of linear equations:
$\frac{2}{x} + \frac{3}{y} = 11$ and $\frac{5}{x} - \frac{4}{y} = -7$

\textbf{Answer:}
\begin{center}
\fitmath{x=1, y=1}
\end{center}

\textbf{Solution:}
\begin{enumerate}
  \item Let $p = \frac{1}{x}$ and $q = \frac{1}{y}$ to convert the equations into linear form.
  \item Substitute to get: $2p + 3q = 11$ and $5p - 4q = -7$.
  \item Solve the linear equations using elimination: Multiply first equation by 4 and second by 3, then add to eliminate $q$: $8p + 12q = 44$ and $15p - 12q = -21$ => $23p = 23$ => $p = 1$.
  \item Substitute $p=1$ into $2p+3q=11$: $2(1)+3q=11$ => $3q=9$ => $q=3$. Thus $x = \frac{1}{p} = 1$ and $y = \frac{1}{q} = \frac{1}{3}$.
\end{enumerate}
\textbf{Derivation:}
Given equations:
 $\frac{2}{x} + \frac{3}{y} = 11$   ...(i)
 $\frac{5}{x} - \frac{4}{y} = -7$  ...(ii)
Let  p = $\frac{1}{x}$ ,  q = $\frac{1}{y}$ .
Then (i) becomes  2p+3q=11 , (ii) becomes  5p-4q=-7 .
Multiply (i) by 4:  8p+12q=44 
Multiply (ii) by 3:  15p-12q=-21 
Adding:  23p=23 $\implies p=1$ 
Substitute  p=1  in  2p+3q=11 :  2+3q=11 $\implies 3q=9 \implies q=3$ 
So  x = $\frac{1}{p} = 1$ ,  y = $\frac{1}{q} = \frac{1}{3}$ .
\hfill \textit{Pair of Linear Equations in Two Variables — equations reducible to linear form}
\vspace{0.3cm}

\question \textbf{20.} A train travels a distance of 480 km at a uniform speed. If the speed had been 8 km/h less, then it would have taken 3 hours more to cover the same distance. Find the usual speed of the train. (Note: Form a quadratic equation and solve it.)

\textbf{Answer:}
\begin{center}
\fitmath{\text{Usual speed of the train} = 40 km/h}
\end{center}

\textbf{Solution:}
\begin{enumerate}
  \item Let the usual speed of the train be $x$ km/h.
  \item Time taken at usual speed = distance/speed = $\frac{480}{x}$ hours.
  \item If speed is 8 km/h less, new speed = $(x - 8)$ km/h.
  \item Time taken at reduced speed = $\frac{480}{x - 8}$ hours.
  \item According to the problem, the time at reduced speed is 3 hours more than the usual time: $\frac{480}{x - 8} = \frac{480}{x} + 3$.
  \item Multiply both sides by $x(x - 8)$ to clear denominators: $480x = 480(x - 8) + 3x(x - 8)$.
  \item Simplify: $480x = 480x - 3840 + 3x^2 - 24x$.
  \item Cancel $480x$ on both sides: $0 = -3840 + 3x^2 - 24x$.
  \item Rearrange: $3x^2 - 24x - 3840 = 0$. Divide through by 3: $x^2 - 8x - 1280 = 0$.
  \item Solve the quadratic equation: $x^2 - 8x - 1280 = 0$. Using quadratic formula $x = \frac{-b \pm \sqrt{b^2 - 4ac}}{2a}$ with $a=1$, $b=-8$, $c=-1280$.
  \item Discriminant $D = (-8)^2 - 4(1)(-1280) = 64 + 5120 = 5184$.
  \item So $x = \frac{8 \pm \sqrt{5184}}{2} = \frac{8 \pm 72}{2}$.
  \item Thus $x = \frac{8 + 72}{2} = 40$ or $x = \frac{8 - 72}{2} = -32$.
  \item Since speed cannot be negative, $x = 40$ km/h.
\end{enumerate}
\textbf{Derivation:}
Let usual speed =  x  km/h. Usual time =  $\frac{480}{x}  h.$ Reduced speed =  (x-8)  km/h. New time =  $\frac{480}{x-8}  h.$ Condition:  $\frac{480}{x-8} = \frac{480}{x} + 3$ . Multiply by  x(x-8) :  480x = 480(x-8) + 3x(x-8) . Simplify:  480x = 480x - 3840 + 3x\textasciicircum{}2 - 24x . Cancel  480x :  0 = 3x\textasciicircum{}2 - 24x - 3840 . Divide by 3:  x\textasciicircum{}2 - 8x - 1280 = 0 . Quadratic formula:  x = $\frac{8 \pm \sqrt{64 + 5120}}{2} = \frac{8 \pm 72}{2}$ . Positive root:  x = 40 .
\hfill \textit{Quadratic Equations — Word Problems (Speed-Distance-Time)}
\vspace{0.3cm}

\question \textbf{21.} In $\triangle ABC$, $P$ and $Q$ are points on sides $AB$ and $AC$ respectively such that $PQ \parallel BC$. The area of $\triangle APQ$ is $16 \text{ cm}^2$ and the area of trapezium $\text{PBCQ}$ is $20 \text{ cm}^2$. Find the ratio $AP:PB$. If $AB = 12$ cm, find the lengths of $AP$ and $PB$.

\textbf{Answer:}
\begin{center}
\fitmath{AP:PB = 2:1, AP = 8 cm, PB = 4 cm}
\end{center}

\textbf{Solution:}
\begin{enumerate}
  \item Let $AP = x$, $PB = y$, so $AB = x+y$.
  \item Since $PQ \parallel BC$, $\triangle APQ \sim \triangle ABC$ (AA similarity).
  \item Ratio of areas of similar triangles equals square of ratio of corresponding sides: $\frac{\text{area}(\triangle APQ)}{\text{area}(\triangle ABC)} = \left(\frac{AP}{AB}\right)^2$.
  \item Area of $\triangle ABC = \text{area}(\triangle APQ) + \text{area}(\text{PBCQ}) = 16 + 20 = 36 \text{ cm}^2$.
  \item Thus $\frac{16}{36} = \left(\frac{x}{x+y}\right)^2 \Rightarrow \frac{4}{9} = \left(\frac{x}{x+y}\right)^2 \Rightarrow \frac{x}{x+y} = \frac{2}{3}$ (taking positive ratio).
  \item Cross-multiply: $3x = 2x + 2y \Rightarrow x = 2y \Rightarrow \frac{x}{y} = \frac{2}{1}$. So $AP:PB = 2:1$.
  \item Given $AB = 12$ cm, so $x + y = 12$. Since $x = 2y$, substitute: $2y + y = 12 \Rightarrow 3y = 12 \Rightarrow y = 4$ cm, then $x = 8$ cm.
\end{enumerate}
\textbf{Derivation:}
\begin{center}
\fitmath{\begin{aligned}
\text{area}(\triangle ABC) &= 16 + 20 = 36 \text{ cm}^2 \\
\frac{\text{area}(\triangle APQ)}{\text{area}(\triangle ABC)} &= \left(\frac{AP}{AB}\right)^2 \\
\frac{16}{36} &= \left(\frac{x}{x+y}\right)^2 \\
\frac{4}{9} &= \left(\frac{x}{x+y}\right)^2 \\
\frac{x}{x+y} &= \frac{2}{3} \\
3x &= 2(x+y) \\
x &= 2y \\
AP:PB &= x:y = 2:1 \\
AB &= x+y = 12 \\
\Rightarrow 2y+y=12 \\
\Rightarrow y=4, x=8
\end{aligned}}
\end{center}
\hfill \textit{Triangles — Similarity and Area Ratio}
\vspace{0.3cm}

\question \textbf{22.} The discriminant of the quadratic equation $x^2 - 5x + 6 = 0$ is:

\textbf{Answer:}
(A) 1

\textbf{Solution:}
\begin{enumerate}
  \item Identify coefficients: a = 1, b = -5, c = 6.
  \item \(
\text{Use discriminant formula D} = b^{2} - 4ac.
\)
  \item Substitute values: D = (-5)$^{2}$ - 4(1)(6) = 25 - 24 = 1.
\end{enumerate}
\textbf{Derivation:}
\begin{center}
\fitmath{D = b^2 - 4ac = (-5)^2 - 4 \cdot 1 \cdot 6 = 25 - 24 = 1}
\end{center}
\hfill \textit{Quadratic Equations — Discriminant}
\vspace{0.3cm}

\question \textbf{23.} The discriminant of the quadratic equation $2x^2 - 3x + 1 = 0$ is:

\textbf{Answer:}
(A) 1

\textbf{Solution:}
\begin{enumerate}
  \item Compare with standard form $ax^2 + bx + c = 0$: $a = 2$, $b = -3$, $c = 1$.
  \item Discriminant formula: $D = b^2 - 4ac$.
  \item Substitute: $D = (-3)^2 - 4(2)(1) = 9 - 8 = 1$.
\end{enumerate}
\textbf{Derivation:}
\begin{center}
\fitmath{D = b^2 - 4ac = (-3)^2 - 4(2)(1) = 9 - 8 = 1}
\end{center}
\hfill \textit{Quadratic Equations — Discriminant and Nature of Roots}
\vspace{0.3cm}

\question \textbf{24.} Assertion (A): The quadratic equation $2x^2 - 5x + 3 = 0$ has two distinct real roots.
Reason (R): For a quadratic equation $ax^2 + bx + c = 0$, if $b^2 - 4ac > 0$, then the equation has two distinct real roots.

\textbf{Answer:}
(A)

\textbf{Solution:}
\begin{enumerate}
  \item Step 1: Identify coefficients: a = 2, b = -5, c = 3.
  \item Step 2: Compute discriminant: D = b$^{2}$ - 4ac = (-5)$^{2}$ - 4(2)(3) = 25 - 24 = 1.
  \item Step 3: Since D = 1 > 0, the quadratic equation has two distinct real roots. Thus Assertion (A) is true.
  \item Step 4: Reason (R) states the correct condition for two distinct real roots, which is true and correctly explains why (A) is true.
\end{enumerate}
\textbf{Derivation:}
\begin{center}
\fitmath{D = b^2 - 4ac = (-5)^2 - 4(2)(3) = 25 - 24 = 1 > 0}
\end{center}
\hfill \textit{Quadratic Equations — Nature of Roots}
\vspace{0.3cm}

\question \textbf{25.} Find the discriminant of the quadratic equation $2x^2 - 4x + 3 = 0$.

\textbf{Answer:}
\begin{center}
\fitmath{-8}
\end{center}

\textbf{Solution:}
\begin{enumerate}
  \item Recall the standard form of a quadratic equation: $ax^2 + bx + c = 0$.
  \item Identify coefficients: $a = 2$, $b = -4$, $c = 3$.
  \item Use the discriminant formula: $D = b^2 - 4ac$.
  \item Substitute values: $D = (-4)^2 - 4(2)(3) = 16 - 24 = -8$.
\end{enumerate}
\textbf{Derivation:}
\begin{center}
\fitmath{D = b^2 - 4ac = (-4)^2 - 4 \cdot 2 \cdot 3 = 16 - 24 = -8}
\end{center}
\hfill \textit{Quadratic Equations — Nature of Roots}
\vspace{0.3cm}

\question \textbf{26.} The discriminant of the quadratic equation $2x^2 - 3x + 5 = 0$ is:

\textbf{Answer:}
(B) $-31$

\textbf{Solution:}
\begin{enumerate}
  \item Identify coefficients: $a = 2$, $b = -3$, $c = 5$
  \item Use discriminant formula: $D = b^2 - 4ac$
  \item Substitute values: $D = (-3)^2 - 4(2)(5) = 9 - 40 = -31$
  \item Thus the discriminant is $-31$.
\end{enumerate}
\textbf{Derivation:}
\begin{center}
\fitmath{D = b^2 - 4ac = (-3)^2 - 4(2)(5) = 9 - 40 = -31}
\end{center}
\hfill \textit{Quadratic Equations — Nature of Roots}
\vspace{0.3cm}

\question \textbf{27.} Find the discriminant of the quadratic equation $2x^2 - 3x + 1 = 0$.

\textbf{Answer:}
\begin{center}
\fitmath{1}
\end{center}

\textbf{Solution:}
\begin{enumerate}
  \item Identify coefficients: $a = 2$, $b = -3$, $c = 1$.
  \item Apply discriminant formula: $D = b^2 - 4ac$.
  \item Compute: $D = (-3)^2 - 4(2)(1) = 9 - 8 = 1$.
\end{enumerate}
\textbf{Derivation:}
\begin{center}
\fitmath{D = (-3)^2 - 4 \cdot 2 \cdot 1 = 9 - 8 = 1}
\end{center}
\hfill \textit{Quadratic Equations — Discriminant}
\vspace{0.3cm}

\question \textbf{28.} Assertion (A): The quadratic equation $x^2 - 3x + 5 = 0$ has no real roots.
Reason (R): For a quadratic equation $ax^2+bx+c=0$, if discriminant $D = b^2 - 4ac < 0$, then the equation has no real roots.

\textbf{Answer:}
(A)

\textbf{Solution:}
\begin{enumerate}
  \item Calculate the discriminant of the quadratic equation $x^2 - 3x + 5 = 0$: $a=1$, $b=-3$, $c=5$.
  \item $D = b^2 - 4ac = (-3)^2 - 4(1)(5) = 9 - 20 = -11$.
  \item Since $D < 0$, the equation has no real roots, making Assertion (A) true.
  \item Reason (R) correctly states the condition for no real roots, and it explains why (A) is true.
\end{enumerate}
\textbf{Derivation:}
\begin{center}
\fitmath{D = (-3)^2 - 4(1)(5) = 9 - 20 = -11 < 0}
\end{center}
\hfill \textit{Quadratic Equations — Nature of roots}
\vspace{0.3cm}

\question \textbf{29.} The roots of the quadratic equation $x^2 - 7x + 10 = 0$ are:

\textbf{Answer:}
(A)

\textbf{Solution:}
\begin{enumerate}
  \item Identify the coefficients: $a = 1$, $b = -7$, $c = 10$.
  \item Calculate the discriminant: $D = b^2 - 4ac = (-7)^2 - 4(1)(10) = 49 - 40 = 9 > 0$.
  \item Since $D > 0$, roots are real and distinct.
  \item Use quadratic formula: $x = \frac{-b \pm \sqrt{D}}{2a} = \frac{7 \pm 3}{2}$.
  \item So $x = \frac{7+3}{2} = 5$ and $x = \frac{7-3}{2} = 2$.
\end{enumerate}
\textbf{Derivation:}
\begin{center}
\fitmath{\begin{aligned}
x = \frac{7 \pm \sqrt{49 - 40}}{2} = \frac{7 \pm 3}{2} \\
\Rightarrow x = 5, 2
\end{aligned}}
\end{center}
\hfill \textit{Quadratic Equations — Quadratic Formula and Nature of Roots}
\vspace{0.3cm}

\question \textbf{30.} Find the discriminant of the quadratic equation $2x^2 - 4x + 3 = 0$.

\textbf{Answer:}
\begin{center}
\fitmath{-8}
\end{center}

\textbf{Solution:}
\begin{enumerate}
  \item Write the quadratic equation in standard form: $2x^2 - 4x + 3 = 0$.
  \item Identify coefficients: $a = 2$, $b = -4$, $c = 3$.
  \item Use discriminant formula: $D = b^2 - 4ac$.
  \item Substitute: $D = (-4)^2 - 4(2)(3) = 16 - 24 = -8$.
\end{enumerate}
\textbf{Derivation:}
\begin{center}
\fitmath{D = b^2 - 4ac = (-4)^2 - 4 \cdot 2 \cdot 3 = 16 - 24 = -8}
\end{center}
\hfill \textit{Quadratic Equations — Nature of Roots}
\vspace{0.3cm}

\question \textbf{31.} Which of the following is the discriminant of the quadratic equation $2x^2 - 4x + 3 = 0$?

\textbf{Answer:}
(C) $16 - 24$

\textbf{Solution:}
\begin{enumerate}
  \item Identify coefficients: $a = 2$, $b = -4$, $c = 3$.
  \item Use discriminant formula: $D = b^2 - 4ac$.
  \item Substitute: $D = (-4)^2 - 4(2)(3) = 16 - 24 = -8$.
\end{enumerate}
\textbf{Derivation:}
\begin{center}
\fitmath{D = b^2 - 4ac = (-4)^2 - 4 \cdot 2 \cdot 3 = 16 - 24 = -8}
\end{center}
\hfill \textit{Quadratic Equations — Nature of roots}
\vspace{0.3cm}

\question \textbf{32.} Find the discriminant of the quadratic equation $2x^2 - 4x + 3 = 0$. Hence, determine the nature of its roots.

\textbf{Answer:}
\begin{center}
\fitmath{\text{Discriminant} = -8; \text{No real roots}.}
\end{center}

\textbf{Solution:}
\begin{enumerate}
  \item Identify coefficients: $a = 2$, $b = -4$, $c = 3$.
  \item Compute discriminant: $D = b^2 - 4ac = (-4)^2 - 4(2)(3) = 16 - 24 = -8$.
  \item Since $D < 0$, the quadratic equation has no real roots.
\end{enumerate}
\textbf{Derivation:}
 D = b\textasciicircum{}2 - 4ac = (-4)\textasciicircum{}2 - 4 $\cdot 2 \cdot 3 = 16 - 24 = -8$ .
 D < 0 $\Rightarrow$  no real roots.
\hfill \textit{Quadratic Equations — Nature of Roots}
\vspace{0.3cm}

\question \textbf{33.} Find the nature of the roots of the quadratic equation $2x^2 - 4x + 3 = 0$.

\textbf{Answer:}
\begin{center}
\fitmath{\text{No real roots} (\text{roots are imaginary})}
\end{center}

\textbf{Solution:}
\begin{enumerate}
  \item Step 1: Identify coefficients $a=2$, $b=-4$, $c=3$ and compute discriminant $D = b^2 - 4ac$.
  \item Step 2: Substitute values: $D = (-4)^2 - 4(2)(3) = 16 - 24 = -8$. Since $D < 0$, the quadratic has no real roots.
\end{enumerate}
\textbf{Derivation:}
\begin{center}
\fitmath{D = b^2 - 4ac = (-4)^2 - 4 \cdot 2 \cdot 3 = 16 - 24 = -8 < 0}
\end{center}
\hfill \textit{Quadratic Equations — Nature of Roots}
\vspace{0.3cm}

\question \textbf{34.} Find the nature of the roots of the quadratic equation $2x^2 - 4x + 3 = 0$.

\textbf{Answer:}
\begin{center}
\fitmath{\text{No real roots}}
\end{center}

\textbf{Solution:}
\begin{enumerate}
  \item Step 1: Identify coefficients: $a = 2$, $b = -4$, $c = 3$.
  \item Step 2: Compute discriminant $D = b^2 - 4ac = (-4)^2 - 4(2)(3) = 16 - 24 = -8$.
  \item Step 3: Since $D < 0$, the quadratic equation has no real roots.
\end{enumerate}
\textbf{Derivation:}
\begin{center}
\fitmath{D = b^2 - 4ac = (-4)^2 - 4 \cdot 2 \cdot 3 = 16 - 24 = -8 < 0}
\end{center}
\hfill \textit{Quadratic Equations — Nature of Roots}
\vspace{0.3cm}

\question \textbf{35.} Find the discriminant of the quadratic equation $3x^2 - 5x + 2 = 0$ and hence determine the nature of its roots.

\textbf{Answer:}
\begin{center}
\fitmath{\text{Discriminant} = 1; \text{Roots are real and distinct}.}
\end{center}

\textbf{Solution:}
\begin{enumerate}
  \item Step 1: Compare the given equation with standard form $ax^2+bx+c=0$ to get $a=3$, $b=-5$, $c=2$. Compute discriminant $D = b^2 - 4ac$.
  \item Step 2: Substitute values: $D = (-5)^2 - 4(3)(2) = 25-24=1$. Since $D>0$, the roots are real and distinct.
\end{enumerate}
\textbf{Derivation:}
\begin{center}
\fitmath{D = b^2-4ac = (-5)^2 - 4(3)(2) = 25-24 = 1 > 0}
\end{center}
\hfill \textit{Quadratic Equations — Nature of roots}
\vspace{0.3cm}

\question \textbf{36.} Find the discriminant of the quadratic equation $3x^2 - 6x + 2 = 0$ and hence determine the nature of its roots.

\textbf{Answer:}
\begin{center}
\fitmath{\text{Discriminant} = 12; \text{Roots are real and distinct}.}
\end{center}

\textbf{Solution:}
\begin{enumerate}
  \item Compare with $ax^2 + bx + c = 0$: $a = 3$, $b = -6$, $c = 2$. Compute discriminant: $D = b^2 - 4ac = (-6)^2 - 4(3)(2) = 36 - 24 = 12$.
  \item Since $D > 0$, the roots are real and distinct.
\end{enumerate}
\textbf{Derivation:}
 D = b\textasciicircum{}2 - 4ac = (-6)\textasciicircum{}2 - 4(3)(2) = 36 - 24 = 12 > 0  $\Rightarrow$  Roots are real and distinct.
\hfill \textit{Quadratic Equations — Nature of roots}
\vspace{0.3cm}

\question \textbf{37.} A train travels 360 km at a uniform speed. If the speed had been 5 km/h more, it would have taken 1 hour less for the same journey. Find the speed of the train.

\textbf{Answer:}
\begin{center}
\fitmath{40 km/h}
\end{center}

\textbf{Solution:}
\begin{enumerate}
  \item Let the original speed of the train be $x$ km/h.
  \item Time taken at original speed = $\frac{360}{x}$ hours. Time taken at increased speed ($x+5$ km/h) = $\frac{360}{x+5}$ hours.
  \item According to the problem, $\frac{360}{x} - \frac{360}{x+5} = 1$. Multiply both sides by $x(x+5)$: $360(x+5) - 360x = x(x+5)$, which simplifies to $1800 = x^2 + 5x$.
  \item Rearrange: $x^2 + 5x - 1800 = 0$. Factorising: $(x+45)(x-40)=0$. Since speed cannot be negative, $x=40$.
\end{enumerate}
\textbf{Derivation:}
\begin{center}
\fitmath{\begin{aligned}
\frac{360}{x} - \frac{360}{x+5} &= 1 \\
360\left(\frac{1}{x} - \frac{1}{x+5}\right) &= 1 \\
360\cdot\frac{5}{x(x+5)} &= 1 \\
1800 &= x^2+5x \\
x^2+5x-1800&=0 \\
(x+45)(x-40)&=0 \\
x=40 \;\text{(taking positive value)}
\end{aligned}}
\end{center}
\hfill \textit{Quadratic Equations — Word problems on speed, distance and time}
\vspace{0.3cm}

\question \textbf{38.} The speed of a boat in still water is 8 km/h. It takes the same time to travel 15 km upstream as it takes to travel 25 km downstream. Find the speed of the stream.

\textbf{Answer:}
\begin{center}
\fitmath{2 km/h}
\end{center}

\textbf{Solution:}
\begin{enumerate}
  \item Step 1: Let the speed of the stream be $x$ km/h. Then, speed upstream $= (8 - x)$ km/h and speed downstream $= (8 + x)$ km/h.
  \item Step 2: Time upstream $= \frac{15}{8-x}$ hours, time downstream $= \frac{25}{8+x}$ hours. Since both times are equal, we have $\frac{15}{8-x} = \frac{25}{8+x}$.
  \item Step 3: Cross-multiplying gives $15(8+x) = 25(8-x)$. Simplify: $120 + 15x = 200 - 25x$. So $15x + 25x = 200 - 120$, i.e. $40x = 80$, hence $x = 2$.
  \item Step 4: Verify: upstream speed = 6 km/h, time = 15/6 = 2.5 h; downstream speed = 10 km/h, time = 25/10 = 2.5 h. Thus the speed of the stream is 2 km/h.
\end{enumerate}
\textbf{Derivation:}
Let speed of stream  = x  km/h. Upstream speed  = 8-x , downstream  = 8+x .
 $\frac{15}{8-x} = \frac{25}{8+x} \implies 15(8+x)=25(8-x) \implies 120+15x=200-25x \implies 40x=80 \implies x=2$ .
\hfill \textit{Quadratic Equations — Word Problems (speed and time)}
\vspace{0.3cm}

\question \textbf{39.} The sum of the squares of two consecutive positive integers is 365. Find the integers.

\textbf{Answer:}
\begin{center}
\fitmath{13 and 14}
\end{center}

\textbf{Solution:}
\begin{enumerate}
  \item Let the two consecutive positive integers be $x$ and $x+1$.
  \item According to the problem: $x^2 + (x+1)^2 = 365$.
  \item Simplify: $x^2 + x^2 + 2x + 1 = 365$ $\Rightarrow$ $2x^2 + 2x - 364 = 0$ $\Rightarrow$ $x^2 + x - 182 = 0$.
  \item Factorize: $(x + 14)(x - 13) = 0$ $\Rightarrow$ $x = 13$ or $x = -14$. Since positive, $x = 13$. Thus the integers are $13$ and $14$.
\end{enumerate}
\textbf{Derivation:}
\begin{center}
\fitmath{\begin{aligned}
&\text{Let integers be }x,\,x+1. \\
&x^2+(x+1)^2=365 \\
\Rightarrow x^2+x^2+2x+1=365 \\
\Rightarrow 2x^2+2x-364=0 \\
\Rightarrow x^2+x-182=0 \\
\Rightarrow (x+14)(x-13)=0 \\
\Rightarrow x=13\quad[\because x>0] \\
&\therefore \text{Integers are }13\text{ and }14.
\end{aligned}}
\end{center}
\hfill \textit{Quadratic Equations — Word Problems}
\vspace{0.3cm}

\question \textbf{40.} The sum of two numbers is 15. If the sum of their reciprocals is $\frac{3}{10}$, find the numbers.

\textbf{Answer:}
\begin{center}
\fitmath{\text{The numbers are} 10 and 5.}
\end{center}

\textbf{Solution:}
\begin{enumerate}
  \item Let the numbers be $x$ and $15-x$. Their reciprocals sum to $\frac{1}{x} + \frac{1}{15-x} = \frac{3}{10}$.
  \item Multiply both sides by $10x(15-x)$: $10(15-x) + 10x = 3x(15-x)$, which simplifies to $150 = 45x - 3x^2$.
  \item Rearrange into standard quadratic: $3x^2 - 45x + 150 = 0$. Divide by 3: $x^2 - 15x + 50 = 0$.
  \item Factorise: $(x-10)(x-5)=0$, so $x=10$ or $x=5$. Thus the numbers are 10 and 5.
\end{enumerate}
\textbf{Derivation:}
Let numbers be  x  and  15-x .
 $\frac{1}{x} + \frac{1}{15-x} = \frac{3}{10}$ 
 10(15-x) + 10x = 3x(15-x) 
 150 = 45x - 3x\textasciicircum{}2 
 3x\textasciicircum{}2 - 45x + 150 = 0 
 x\textasciicircum{}2 - 15x + 50 = 0 
 (x-10)(x-5)=0 
 x=10  or  x=5 
\hfill \textit{Quadratic Equations — Word problems}
\vspace{0.3cm}

\end{questions}

\end{document}
